% !TEX root = ../main.tex

\begin{table}[htbp]
  \centering
  \small
  \caption[Evaluación del sistema completo]{%
    Comportamiento del sistema completo sobre las 57
      entradas del banco de evaluación, por familia. Fuente: elaboración
      propia. \textcolor{borrador2}{Medición del 06/09/2026 (IT-124).}}
  \label{tab:evaluacion_sistema}
  \begin{tabular}{@{}lrrr@{}}
    \toprule
    \textbf{Familia}                    & \textbf{n} & \textbf{Aciertos}
                                        & \textbf{Mediana (s)}             \\
    \midrule
    Créditos de una asignatura          & 4          & 4  & \textcolor{borrador2}{18,2} \\
    Curso en que se imparte             & 4          & 4  & \textcolor{borrador2}{25,9} \\
    Asignaturas de una mención          & 4          & 4  & \textcolor{borrador2}{34,4} \\
    Plan de estudios de un curso        & 4          & 4  & \textcolor{borrador2}{40,4} \\
    Optativas de una titulación         & 4          & 4  & \textcolor{borrador2}{62,5} \\
    Catálogo del centro                 & 1          & 1  & \textcolor{borrador2}{67,4} \\
    \midrule
    Conversación encadenada             & 8          & 8  & \textcolor{borrador2}{81,9} \\
    Petición de consejo                 & 6          & 6  & \textcolor{borrador2}{98,3} \\
    Pregunta ambigua                    & 3          & 3  & \textcolor{borrador2}{124,3} \\
    Cortesía                            & 4          & 4  & \textcolor{borrador2}{0,0} \\
    Fuera de dominio                    & 15         & 15 & \textcolor{borrador2}{0,1} \\
    \midrule
    \textbf{Total}                      & \textbf{57} & \textbf{57} & \textcolor{borrador2}{\textbf{39,4}} \\
    \bottomrule
  \end{tabular}
  \begin{minipage}{0.95\textwidth}
    \vspace{4pt}
    \tiny
    \textbf{Notas:} $n$ es el número de casos de prueba evaluados. Todas las comprobaciones se corrigen de forma determinista contra los datos de la colección o las respuestas fijas del sistema (sin modelos como juez). En \textit{Conversación encadenada}, el tiempo acumula la totalidad de turnos del diálogo. Las familias de \textit{Cortesía} y \textit{Fuera de dominio} presentan latencias mínimas porque la mayoría de sus entradas se resuelven mediante respuestas fijas o filtros antes de invocar al modelo generativo.
  \end{minipage}
\end{table}
